\mxunit{premiere-c08}{Variables aléatoires et dispersion}\label{C08}
\begin{keypoint}[Le parcours du chapitre]
\textbf{Objectifs.} Construire une loi ; interpréter des événements ; calculer espérance, variance, écart type ; utiliser la linéarité ; simuler des échantillons et calculer leur moyenne observée.

\textbf{Prérequis.} Probabilités, moyenne pondérée, carrés, arbres et listes.

\textbf{Repère annuel.} 19–28 janvier : 8 h\end{keypoint}
\section*{Diagnostic : \mxcount{diag} questions, 10 minutes}\label{C08-D}
Répondre sans consulter le cours. Indiquer les questions pour lesquelles une aide est nécessaire. Ce diagnostic ne donne pas lieu à une note.
\begin{enumerate}
\item Calculer la moyenne de 2, 4, 6.
\item Vérifier si $\pt{\dfrac{1}{2};\dfrac{1}{3};\dfrac{1}{6}}$ peut être une loi.
\item Calculer $(2-5)^2$.
\item Si $P(A)=\num{0,6}$, calculer $P(\overline A)$.
\item Un gain brut de 8 unités et une mise de 3 donnent quel gain net ?
\item Calculer $\sqrt{\num{0,25}}$.
\mxcounted{diag}
\end{enumerate}
\paragraph{Après le diagnostic.} Reprendre les notions signalées avec le professeur ; refaire une question analogue, puis expliquer la règle utilisée.
\paragraph{Oral et logique.} Expliquer pourquoi une moyenne théorique peut être un gain impossible. Présenter la différence entre une loi, une simulation et une série statistique.
\mxnewpage
\section*{Chercher avant de formaliser}\label{C08-RECH}
Créer deux jeux à trois issues équiprobables ayant espérance 0 mais des écarts types distincts. Vérifier leurs lois et expliquer la différence.
\paragraph{Organisation.} Recherche individuelle 5 min, confrontation à deux 5 min, mise en commun 5 min. Garder une trace d’un essai, même non abouti. Une conjecture doit ensuite être justifiée.
\mxnewpage
\section{Variable, événement et loi}\label{C08-U1}
\begin{keypoint}[Définition et résultats]
Sur un univers fini $\Omega$, une variable aléatoire réelle X est une fonction qui associe un nombre réel à chaque issue. L’événement $\{X=a\}$ regroupe les issues donnant la valeur a. Sa loi associe à chaque valeur distincte $x_i$ la probabilité $p_i=P(X=x_i)$, avec $p_i\geq0$ et $\sum p_i=1$.
\end{keypoint}
\paragraph{Statut et justification.} Définitions. Les événements associés aux valeurs distinctes forment une partition de l’univers ; cela justifie la somme des probabilités.
\begin{example}Un dé équilibré donne un gain X égal à 2 pour 1 ou 2, et à $-1$ sinon. La loi est $P(X=2)=\dfrac{1}{3}$, $P(X=-1)=\dfrac{2}{3}$. $\{X\leq0\}=\ens{3;4;5;6}$ : probabilité $\dfrac{2}{3}$.\end{example}\label{C08-EX1}
\mxinclude{figures/C08-U1}
\begin{pitfall}Les valeurs ne sont pas toujours équiprobables même si les issues le sont. Regrouper les issues qui donnent le même gain.\end{pitfall}
\begin{application}Deux pièces équilibrées sont lancées. X compte les piles. Écrire sa loi et $P(X\leq1)$.\end{application}\label{C08-AP1}
\section{Espérance et linéarité}\label{C08-U2}
\begin{keypoint}[Définition et résultats]
L’espérance est $E(X)=\sum_i x_ip_i$. Elle décrit la moyenne théorique à long terme, pas un résultat nécessairement possible sur un essai. Pour des variables X et Y sur le même univers fini et des réels a,b : $E(aX+b)=aE(X)+b$ et $E(X+Y)=E(X)+E(Y)$, sans hypothèse d’indépendance.
\end{keypoint}
\paragraph{Statut et justification.} En sommant sur les issues $\omega$, $E(X+Y)=\sum_\omega(X(\omega)+Y(\omega))P(\omega)$ ; distribuer la somme donne le résultat. La constante b a espérance b car les probabilités ont somme 1.
\begin{example}Pour la loi précédente $P(X=2)=\dfrac{1}{3}$, $P(X=-1)=\dfrac{2}{3}$, $E(X)=0$ : le jeu est équitable si X est le gain net. Pour $Y=3X-2$, $E(Y)=3\times0-2=-2$.\end{example}\label{C08-EX2}
\begin{pitfall}Une espérance n’est pas une garantie individuelle ; un jeu peut être perdant à chaque horizon observé malgré une espérance positive.\end{pitfall}
\begin{application}Un gain brut vaut 10 avec probabilité $\dfrac{1}{5}$ et 0 sinon. Quelle mise rend le jeu équitable ?\end{application}\label{C08-AP2}
\section{Variance, écart type et König–Huygens}\label{C08-U3}
\begin{keypoint}[Définition et résultats]
En notant $\mu=E(X)$, $V(X)=E((X-\mu)^2)=\sum_i(x_i-\mu)^2p_i$. L’écart type est $\sigma(X)=\sqrt{V(X)}$. La formule de König–Huygens donne $V(X)=E(X^2)-E(X)^2$. Une variance est toujours positive ou nulle ; l’écart type a la même unité que X, la variance une unité au carré.
\end{keypoint}
\paragraph{Statut et justification.} Développer $(X-\mu)^2=X^2-2\mu X+\mu^2$, appliquer la linéarité et $E(X)=\mu$ : $E(X^2)-2\mu^2+\mu^2$. On obtient la formule.
\begin{example}Pour X de valeurs 0 et 4 avec probabilités $\dfrac{3}{4}$ et $\dfrac{1}{4}$, $E(X)=1$, $E(X^2)=4$, $V(X)=3$ et $\sigma=\sqrt3$. Calcul direct : $(0-1)^2(\dfrac{3}{4})+(4-1)^2(\dfrac{1}{4})=3$.\end{example}\label{C08-EX3}
\begin{pitfall}$E(X^2)$ n’est pas $E(X)^2$. Une variance calculée négative signale une erreur ou une très petite erreur d’arrondi à traiter explicitement.\end{pitfall}
\begin{application}Pour X de valeurs $-1,1$ équiprobables, calculer espérance, variance et écart type.\end{application}\label{C08-AP3}
\section{Transformer et comparer des gains}\label{C08-U4}
\begin{keypoint}[Définition et résultats]
Pour a,b réels, $V(aX+b)=a^2V(X)$ et $\sigma(aX+b)=|a|\sigma(X)$. Une translation déplace la moyenne sans changer la dispersion ; une multiplication change l’échelle. Deux jeux de même espérance peuvent avoir des dispersions très différentes.
\end{keypoint}
\paragraph{Statut et justification.} $aX+b-E(aX+b)=a(X-E(X))$. En élevant au carré puis en prenant l’espérance, on obtient $a^2V(X)$. La racine carrée de $a^2$ est $|a|$.
\begin{example}Un jeu A rapporte toujours 2. Un jeu B rapporte 0 ou 4 équiprobablement. Les deux ont espérance 2 ; variances 0 et 4. Leurs distributions ne sont donc pas interchangeables malgré la même moyenne.\end{example}\label{C08-EX4}
\begin{pitfall}La variance d’une somme n’est pas, en général, la somme des variances ; aucune règle de ce type n’est utilisée sans hypothèses.\end{pitfall}
\begin{application}$E(X)=3$, $V(X)=4$. Donner les caractéristiques de $Y=-2X+5$.\end{application}\label{C08-AP4}
\section{Échantillons et moyenne observée}\label{C08-U5}
\begin{keypoint}[Définition et résultats]
Un échantillon de taille n est ici une liste de n réalisations indépendantes d’une même variable X. Sa moyenne observée m varie d’un échantillon à l’autre ; on la compare à $\mu=E(X)$.
\end{keypoint}
\paragraph{Statut et justification.} Expérimentation : on ne démontre ici ni théorème de convergence ni loi normale. On vérifie la taille des listes, le calcul de la moyenne et la distinction entre n et le nombre N d’échantillons.
\begin{example}Pour X de valeurs 0 et 2 équiprobables, $\mu=1$. L’échantillon 2, 0, 2, 2, 0, 0, 2, 2, 0, 2 a pour moyenne observée \num{1,2} ; un autre échantillon de même taille donnera en général une autre moyenne. Une simulation compare plusieurs moyennes observées à $\mu$ sans rien démontrer.\end{example}\label{C08-EX5}
\begin{pitfall}Ne pas recopier un seul tirage n fois ; ne pas confondre la moyenne observée d’un échantillon avec l’espérance $\mu$.\end{pitfall}
\begin{beyond}[Seuil de fluctuation]En répétant la simulation, on peut compter la proportion des échantillons vérifiant $|m-\mu|\leq \dfrac{2\sigma}{\sqrt n}$ ; elle est souvent voisine de \pc{95} sans l’être nécessairement. Ce seuil ne figure pas au programme de première et n’est pas évalué.\end{beyond}
\begin{application}X prend les valeurs 0 et 2 de façon équiprobable. Calculer la moyenne observée de l’échantillon 2, 0, 2, 2, 0, puis son écart à $E(X)$.\end{application}\label{C08-AP5}
\mxnewpage
\section*{Parcours d’exercices}\label{C08-PARC}
Commencer par les exercices 1 et 2. Choisir ensuite un exercice 3 ou 4 sans aide. Les exercices 5 à 7 sont une réserve de transfert, de recherche ou de programmation : le professeur en désigne un. L’exercice 8 se reprend à distance.
\begingroup
\mxexercisehead{C08-E1}{Une loi à construire}
On lance un dé équilibré à six faces numérotées de 1 à 6. La variable aléatoire $X$ associe au résultat le reste de sa division euclidienne par 3. Par exemple, un résultat divisible par 3 donne $X=0$.
\begin{enumerate}[label=\arabic*.]
\item Pour chacune des six faces, indiquer la valeur prise par $X$.
\item Présenter la loi de probabilité de $X$ dans un tableau et vérifier que la somme des probabilités vaut 1.
\item Calculer $P(X<2)$ en précisant quelles valeurs de $X$ sont retenues.
\end{enumerate}
\mxexercisehead{C08-E2}{Deux calculs de variance}
La variable aléatoire $X$ prend les valeurs $-2$, 0 et 4 avec les probabilités respectives $\dfrac14$, $\dfrac12$ et $\dfrac14$.
\begin{enumerate}[label=\arabic*.]
\item Calculer son espérance $\mu=E(X)$ et le moment $E(X^2)$.
\item Calculer $V(X)$ à l’aide de la formule de König–Huygens, puis en déduire l’écart type exact.
\item Retrouver la variance à partir de la définition $V(X)=E((X-\mu)^2)$. Présenter les trois contributions au calcul.
\end{enumerate}
\mxexercisehead{C08-E3}{Un jeu équitable}
Une roue comporte cinq secteurs de même angle, portant les nombres 0, 0, 0, 5 et 10. On suppose les cinq secteurs équiprobables. Le nombre obtenu représente un gain brut $G$, en unités monétaires. Pour jouer, on paie une mise $m$ ; le gain net est donc $X=G-m$.
\begin{enumerate}[label=\arabic*.]
\item Établir la loi de $G$, puis déterminer la mise $m$ qui rend le jeu équitable, c’est-à-dire telle que $E(X)=0$.
\item Après l’étude de la variance, calculer l’écart type du gain brut. En déduire celui du gain net lorsque la mise est équitable.
\item Expliquer pourquoi modifier la mise change l’espérance du gain net mais pas son écart type.
\end{enumerate}
\mxinclude{figures/C08-E3}
\mxexercisehead{C08-E4}{Même moyenne, autre risque}
Deux placements fictifs sont proposés. Le premier procure avec certitude un gain $X=1$. Le second procure un gain $Y=-2$ avec probabilité $\dfrac12$ et un gain $Y=4$ avec probabilité $\dfrac12$. Les gains sont exprimés dans la même unité monétaire.
\begin{enumerate}[label=\arabic*.]
\item Comparer les espérances des deux gains, puis leurs variances.
\item Une personne déclare : « Ces placements sont équivalents puisqu’ils ont la même moyenne. » Expliquer ce que cette phrase omet et ce que la variance permet ici de distinguer.
\end{enumerate}
\mxexercisehead{C08-E5}{Linéarité sans indépendance}
On lance un dé équilibré numéroté de 1 à 6. On note $X$ le résultat du lancer et l’on définit $Y=7-X$. Les deux variables sont donc obtenues à partir du même lancer.
\begin{enumerate}[label=\arabic*.]
\item Calculer $E(X)$ et $E(Y)$, puis $E(X+Y)$. Vérifier la linéarité de l’espérance sur cet exemple.
\item Les variables $X$ et $Y$ sont-elles indépendantes ? Justifier en utilisant, par exemple, les événements $\{X=1\}$ et $\{Y=6\}$.
\item La linéarité de l’espérance exige-t-elle l’indépendance ? Conclure en s’appuyant sur vos résultats.
\end{enumerate}
\mxexercisehead{C08-E6}{Une perte minimale}
Une variable aléatoire $X$ prend un nombre fini de valeurs et possède une espérance $\mu$. Pour un réel $a$, on mesure l’erreur quadratique moyenne associée à la prévision constante $a$ par $E((X-a)^2)$.
\begin{enumerate}[label=\arabic*.]
\item Démontrer que $E((X-a)^2)=V(X)+(a-\mu)^2$. On pourra écrire $X-a=(X-\mu)+(\mu-a)$.
\item En déduire la valeur de $a$ qui minimise cette erreur, puis la valeur minimale. Justifier l’unicité du choix optimal.
\end{enumerate}
\mxexercisehead{C08-E7}{Fréquences de lettres}
On souhaite mesurer la fréquence de la lettre « a » dans un texte. Le dénominateur est le nombre de caractères alphabétiques : les espaces, chiffres et signes de ponctuation sont exclus. Les majuscules sont assimilées aux minuscules ; les lettres accentuées restent distinctes de « a ».
\begin{enumerate}[label=\arabic*.]
\item Écrire une fonction Python qui renvoie cette fréquence. Préciser le résultat choisi lorsque le texte ne contient aucune lettre.
\item Tester la fonction avec \code{Aba !}, puis avec \code{123 !} et avec la chaîne vide. Expliquer les résultats attendus avant d’exécuter le programme.
\end{enumerate}
\mxexercisehead{C08-E8}{Reprise différée}
On répète deux fois, de façon indépendante, une expérience dont la probabilité de succès vaut $\dfrac12$. La variable aléatoire $X$ compte le nombre de succès obtenus.
\begin{enumerate}[label=\arabic*.]
\item Retrouver la loi de $X$ à l’aide des quatre suites de résultats possibles.
\item Calculer $E(X)$ et $V(X)$ en partant de cette loi.
\end{enumerate}
\paragraph{Travail personnel.} Effectuer uniquement la tâche désignée à l’issue de la séance, dans la progression datée : 10–20 min maximum. Les exercices dépendent des notions déjà étudiées ; certaines questions sont donc différées. Reprendre l’exercice 8 une semaine après le chapitre (10 min). Si le temps prévu est dépassé, noter la difficulté et apporter les essais ; ne pas accumuler toute la réserve.
\endgroup
\mxneed{12}\section*{À faire dans le manuel CQFD}
Ces exercices se traitent dans votre manuel. La progression indique la séance ; ils remplacent le travail personnel prévu ce jour-là (10 à 20 minutes au total).
\paragraph{CQFD, exercice 28 p. 306, questions toutes.}\label{C08-M1} Calculer l’espérance à partir du tableau de loi ; conserver les probabilités associées à chaque valeur.
\paragraph{CQFD, exercice 31 p. 306, questions toutes ; utiliser la loi de l’exercice 28.}\label{C08-M2} Calculer la variance puis l’écart type, en respectant l’arrondi au dixième demandé par le manuel.
\begingroup\small
\mxnewpage\section*{QCM 1 : comprendre et choisir}\label{C08-Q1}
\mxCount{qcm1} questions pédagogiques à réponse unique. Choisir une réponse par question ; expliquer ensuite un distracteur. Ce format de chapitre ne reproduit pas le format officiel de l’EAM.
\begin{enumerate}
\item Une loi de probabilité a une somme des probabilités égale à
\begin{enumerate}[label=\Alph*.,itemsep=2pt]
\item $0$
\item $1$
\item $n$
\item la moyenne
\end{enumerate}
\item X vaut 0 et 4 avec probabilités $\dfrac{3}{4}$ et $\dfrac{1}{4}$ ; $E(X)=$
\begin{enumerate}[label=\Alph*.,itemsep=2pt]
\item $1$
\item $2$
\item $3$
\item $4$
\end{enumerate}
\item $V(X)$ vaut
\begin{enumerate}[label=\Alph*.,itemsep=2pt]
\item $E(X^2)$
\item $E(X)^2$
\item $E(X^2)+E(X)^2$
\item $E(X^2)-E(X)^2$
\end{enumerate}
\item Si $V(X)=9$, $\sigma(X)=$
\begin{enumerate}[label=\Alph*.,itemsep=2pt]
\item $81$
\item $9$
\item $3$
\item $\dfrac{1}{3}$
\end{enumerate}
\item $E(2X+1)$ vaut
\begin{enumerate}[label=\Alph*.,itemsep=2pt]
\item $2E(X)$
\item $2E(X)+1$
\item $E(X)+1$
\item $E(X)^2+1$
\end{enumerate}
\mxcounted{qcm1}
\end{enumerate}
\endgroup
\begingroup\small
\mxneed{18}\section*{QCM 2 : comprendre et choisir}\label{C08-Q2}
\mxCount{qcm2} questions pédagogiques à réponse unique. Choisir une réponse par question ; expliquer ensuite un distracteur. Ce format de chapitre ne reproduit pas le format officiel de l’EAM.
\begin{enumerate}
\item $V(X+5)$ vaut
\begin{enumerate}[label=\Alph*.,itemsep=2pt]
\item $V(X)+5$
\item $V(X)+25$
\item $V(X)$
\item $5V(X)$
\end{enumerate}
\item $\sigma(-3X)$ vaut
\begin{enumerate}[label=\Alph*.,itemsep=2pt]
\item $3\sigma(X)$
\item $-3\sigma(X)$
\item $9\sigma(X)$
\item $\sigma(X)$
\end{enumerate}
\item L’espérance est nécessairement
\begin{enumerate}[label=\Alph*.,itemsep=2pt]
\item une valeur possible
\item un entier
\item positive
\item aucune de ces trois affirmations
\end{enumerate}
\item $E(X+Y)=E(X)+E(Y)$ exige
\begin{enumerate}[label=\Alph*.,itemsep=2pt]
\item l’indépendance
\item aucune indépendance
\item X=Y
\item X positive
\end{enumerate}
\item Si $E(X)=2$ et $E(X^2)=5$, alors $V(X)=$
\begin{enumerate}[label=\Alph*.,itemsep=2pt]
\item $3$
\item $9$
\item $1$
\item $7$
\end{enumerate}
\mxcounted{qcm2}
\end{enumerate}
\subsection*{Vrai ou faux ? Justifier.}\label{C08-VF}
\begin{enumerate}
\item Une variance peut être négative.
\item Deux variables de même espérance ont la même loi.
\item Ajouter une constante à une variable ne change pas son écart type.
\end{enumerate}
\begin{keypoint}[Bilan personnel]Je construis une loi complète ; je distingue brut et net ; je contrôle les deux formules de variance ; je ne confonds pas fréquence et probabilité.

Pour chaque objectif, indiquer : acquis sans aide ; réussi avec aide ; à retravailler. Écrire un exemple et une erreur à éviter.\end{keypoint}
\endgroup
\ifmxteacher\mxinclude{corriges/C08}\fi
